% bibliografia.tex: bibliografía única del reporte (TT1 + bibitems_nuevos + bibitems_noche, fusionados el 2026-10-18).
% La incluye Technical_Report.tex; los scripts de Docs/capitulos/ la leen de aquí.
\begin{thebibliography}{99}
% --- Referencia 1: OMS ---
  % SIN CITAR (auditoría A1/A4, 2026-10-19): comentada para que no se imprima; descomentar si se cita.
%   \bibitem{WHO2025}
%     World Health Organization. (2025, August 14). \textit{World Health Organization}.
%     Retrieved from WHO Web site:
%     \url{https://www.who.int/news-room/fact-sheets/detail/breast-cancer}

  % --- Referencia 2: GLOBOCAN Data --- 
  \bibitem{GLOBOCAN2022}
    Freihat, O., Sipos, D., \& Kovacs, A. (2025). Global burden and projections of
    breast cancer incidence and mortality to 2050: a comprehensive analysis of
    GLOBOCAN data. \textit{Frontiers in Public Health}, 13, 1622954.
    \url{https://doi.org/10.3389/fpubh.2025.1622954}

  % --- Referencia 3: BCRF Statistics ---
  % SIN CITAR (auditoría A1/A4, 2026-10-19): comentada para que no se imprima; descomentar si se cita.
%   \bibitem{BCRF2025}
%     Breast Cancer Research Foundation. (2025). Breast cancer statistics and resources.
%     Retrieved from \url{https://www.bcrf.org/breast-cancer-statistics-and-resources/}

  % --- Referencia 4: Stop Breast Cancer ---
  % SIN CITAR (auditoría A1/A4, 2026-10-19): comentada para que no se imprima; descomentar si se cita.
%   \bibitem{StopBreastCancer2025}
%     Stop Breast Cancer. (2025). Facts \& figures. Retrieved from
%     \url{https://www.stopbreastcancer.org/information-center/facts-figures/}

  % --- Referencia 5: Global Cancer Statistics ---
  \bibitem{GlobalCancerStatistics2024}
    American Cancer Society. (2024). Global cancer statistics. Retrieved from
    \url{https://pressroom.cancer.org/GlobalCancerStatistics2024}

  % --- Referencia 6: Cáncer de Mama en México ---
  % SIN CITAR (auditoría A1/A4, 2026-10-19): comentada para que no se imprima; descomentar si se cita.
%   \bibitem{CancerMexico2025}
%     Consultores Salud. (2025). Cáncer de mama en México. Retrieved from
%     \url{https://consultorsalud.com.mx/cancer-de-mama-en-mexico/}

  % --- Referencia 7: INEGI Cáncer de Mama ---
  \bibitem{INEGI2025}
    Instituto Nacional de Estadística y Geografía. (2025). Estadísticas de cáncer de mama.
    Retrieved from \url{https://www.inegi.org.mx/contenidos/saladeprensa/aproposito/2025/EAP_CancerMama_25.pdf}

  % --- Referencia 8: National Breast Cancer Foundation - Early Detection ---
  \bibitem{NBCF2026_EarlyDetection}
    National Breast Cancer Foundation. (2026). Early detection of breast cancer.
    Retrieved from \url{https://www.nationalbreastcancer.org/early-detection-of-breast-cancer/}

  % --- Referencia 9: National Breast Cancer Foundation - Breast Cancer Facts ---
  \bibitem{NBCF2026_Facts}
    National Breast Cancer Foundation. (2026). Breast cancer facts \& statistics.
    Retrieved from \url{https://www.nationalbreastcancer.org/breast-cancer-facts/}

  % --- Referencia 10: Windsong Radiology - Early Detection and Survival Rates ---
  % SIN CITAR (auditoría A1/A4, 2026-10-19): comentada para que no se imprima; descomentar si se cita.
%   \bibitem{Windsong2026_EarlyDetection}
%     Windsong Radiology. (2026). How early detection affects breast cancer survival rates.
%     Retrieved from \url{https://windsongwny.com/breast-health/how-early-detection-affects-breast-cancer-survival-rates/}

  % --- Referencia 11: INEGI - Estadísticas de Cáncer de Mama 2023 ---
  \bibitem{INEGI2023}
    Instituto Nacional de Estadística y Geografía. (2023). Estadísticas de cáncer de mama.
    Retrieved from \url{https://www.inegi.org.mx/contenidos/saladeprensa/aproposito/2023/EAP_CMAMA23.pdf}

  % --- Referencia 12: American Cancer Society - Global Cancer Statistics 2024 ---
  % RETIRADA 2026-10-19 (auditoría A1): duplicado de GlobalCancerStatistics2024, sin citar
%   \bibitem{GlobalCancerStatistics2024_2}
%     American Cancer Society. (2024, April). Global cancer statistics 2024. Retrieved from
%     \url{https://pressroom.cancer.org/GlobalCancerStatistics2024}

  % --- Referencia 13: Deep Learning Models on Mammogram Classification ---
  \bibitem{Wang2020DeepLearning}
    Wang, X., Liang, G., Zhang, Y., Blanton, H., Bessinger, Z., \& Jacobs, N. (2020).
    Inconsistent performance of deep learning models on mammogram classification.
    \textit{Journal of the American College of Radiology}, 17(6), 796-803.
    \url{https://doi.org/10.1016/j.jacr.2020.01.006}

  % --- Referencia 14: Global Patterns and Trends in Breast Cancer ---
  \bibitem{Kim2025GlobalBreastCancer}
    Kim, J., Harper, A., McCormack, V., Sung, H., Houssami, N., Morgan, E., Mutebi, M.,
    et al. (2025). Global patterns and trends in breast cancer incidence and mortality
    across 185 countries. \textit{Nature Medicine}, 31(4), 1154-1162.
    \url{https://doi.org/10.1038/s41591-025-03502-3}

  % --- Referencia 15: Quantum Classical Hybrid CNN for Breast Cancer ---
  \bibitem{Xiang2024QuantumCNN}
    Xiang, Q., Li, D., Hu, Z., Yuan, Y., Sun, Y., Zhu, Y., Fu, Y., Jiang, Y., et al. (2024).
    Quantum classical hybrid convolutional neural networks for breast cancer diagnosis.
    \textit{Scientific Reports}, 14(1), 24699.
    \url{https://doi.org/10.1038/s41598-024-74778-7}

  % --- Referencia 16: Quantum Transfer Learning for Breast Cancer Detection ---
  \bibitem{Azevedo2022QuantumTransfer}
    Azevedo, V., Silva, C., \& Dutra, I. (2022). Quantum transfer learning for breast
    cancer detection. \textit{Quantum Machine Intelligence}, 4(1), 5.
    \url{https://doi.org/10.1007/s42484-022-00062-4}

  % --- Referencia 17: Quantum Kernel Methods for Marketing Analytics ---
  % SIN CITAR (auditoría A1/A4, 2026-10-19): comentada para que no se imprima; descomentar si se cita.
%   \bibitem{SaezOrtuño2026QuantumKernel}
%     Sáez Ortuño, L., Forgas Coll, S., \& Ferrara, M. (2026). Quantum kernel methods
%     for marketing analytics with convergence theory and separation bounds.
%     \textit{Scientific Reports}, 16(1), 6645.
%     \url{https://doi.org/10.1038/s41598-026-35793-y}

  % --- Referencia 18: Benchmarking Quantum Machine Learning Kernel Training ---
  \bibitem{AlvarezEstevez2024QuantumKernelBench}
    Alvarez-Estevez, D. (2024). Benchmarking quantum machine learning kernel training
    for classification tasks. arXiv preprint arXiv:2408.10274.
    \url{https://doi.org/10.48550/arXiv.2408.10274}

  % --- Referencia 19: Artificial Intelligence in Medical Imaging ---
  \bibitem{TandFOnline2021}
    % Corregido el 2026-10-19 (auditoría A4-2): era una plantilla sin rellenar; metadatos de insumos/bibitems_nuevos.tex.
    Zebari, D. A., Ibrahim, D. A., Zeebaree, D. Q., Haron, H., Salih, M. S., Dama\v{s}evi\v{c}ius, R., \& Mohammed, M. A. (2021).
    Systematic review of computing approaches for breast cancer detection based computer aided diagnosis using mammogram images.
    \textit{Applied Artificial Intelligence}, 35(15), 2157--2203.
    \url{https://doi.org/10.1080/08839514.2021.2001177}

  % --- Referencia 20: Machine Learning and Deep Learning in Mammography ---
  % SIN CITAR (auditoría A1/A4, 2026-10-19): comentada para que no se imprima; descomentar si se cita.
%   \bibitem{Gao2023MammographyReview}
%     Gao, Y., Lin, J., Zhou, Y., \& Lin, R. (2023). The application of traditional
%     machine learning and deep learning techniques in mammography: a review.
%     \textit{Frontiers in Oncology}, 13, 1213045.
%     \url{https://doi.org/10.3389/fonc.2023.1213045}

  % --- Referencia 21: AI Techniques for Breast Cancer Detection - Multimodal Review ---
  \bibitem{Mashekova2025AIBreastCancer}
    Mashekova, A., Zhao, M. Y., Zarikas, V., Mukhmetov, O., Aidossov, N., Ng, E. Y. K., Wei, D., \& Shapatova, M. (2025).
    Review of artificial intelligence techniques for breast cancer detection with different modalities:
    Mammography, ultrasound, and thermography images. \textit{Bioengineering}, 12(10), 1110.
    \url{https://doi.org/10.3390/bioengineering12101110}

  % --- Referencia 22: Medical Image Enhancement for Breast Cancer Detection ---
  \bibitem{Avci2023MammographyEnhancement}
    Avcı, H., \& Karakaya, J. (2023). A novel medical image enhancement algorithm for breast cancer
    detection on mammography images using machine learning. \textit{Diagnostics}, 13(3), 348.
    \url{https://doi.org/10.3390/diagnostics13030348}

  % --- Referencia 23: AI Pipeline for Mammography Breast Cancer Detection - PRISMA Review ---
  \bibitem{Anez2025AIPipeline}
    Añez, D., Conti, G., Uriarte, J. J., Serrano-Olmedo, J.-J., Martínez-Murillo, R., \& Casanova-Carvajal, O. (2025).
    Artificial intelligence pipeline for mammography-based breast cancer detection: An integrated systematic review and
    large-scale experimental validation. \textit{Medicina}, 61(12), 2237.
    \url{https://doi.org/10.3390/medicina61122237}

  % --- Referencia 24: Deep Convolutional Neural Networks for Mammography ---
  \bibitem{Abdelhafiz2019DeepCNN}
    Abdelhafiz, D., Yang, C., Ammar, R., et al. (2019). Deep convolutional neural networks for mammography:
    advances, challenges and applications. \textit{BMC Bioinformatics}, 20(Suppl 11), 281.
    \url{https://doi.org/10.1186/s12859-019-2823-4}

  % --- Referencia 25: Deep Learning to Improve Breast Cancer Detection on Screening Mammography ---
  \bibitem{Shen2019DeepLearning}
    Shen, L., Margolies, L. R., Rothstein, J. H., et al. (2019). Deep learning to improve breast cancer
    detection on screening mammography. \textit{Scientific Reports}, 9, 12495.
    \url{https://doi.org/10.1038/s41598-019-48995-4}

  % --- Referencia 26: Integrated Framework for Breast Mass Classification Using Residual Neural Networks ---
  \bibitem{Baccouche2022ResidualNN}
    Baccouche, A., Garcia-Zapirain, B., \& Elmaghraby, A. S. (2022). An integrated framework for breast mass
    classification and diagnosis using stacked ensemble of residual neural networks. \textit{Scientific Reports}, 12, 12259.
    \url{https://doi.org/10.1038/s41598-022-15632-6}

  % --- Referencia 27: Deep Learning Classification of Breast Microcalcifications ---
  \bibitem{Shia2025MicrocalcClassification}
    Shia, W. C. (2025). Deep learning-based classification of benign and malignant breast microcalcifications in
    mammography. \textit{Scientific Reports}, 15, 42768.
    \url{https://doi.org/10.1038/s41598-025-27149-9}

  % --- Referencia 28: Quantum Machine Learning in Feature Hilbert Spaces ---
  \bibitem{Schuld2019QuantumML}
    Schuld, M., \& Killoran, N. (2019). Quantum machine learning in feature {H}ilbert spaces.
    \textit{Physical Review Letters}, 122(4), 040504.
    \url{https://doi.org/10.1103/PhysRevLett.122.040504}

  % --- Referencia 29: Hybrid Quantum-Classical-Quantum Convolutional Neural Networks ---
  \bibitem{Long2025HybridQCQ}
    Long, C., Huang, M., Ye, X., et al. (2025). Hybrid quantum-classical-quantum convolutional neural networks.
    \textit{Scientific Reports}, 15, 31780.
    \url{https://doi.org/10.1038/s41598-025-13417-1}

  % --- Referencia 30: Quantum Machine and Deep Learning for Medical Image Classification ---
  % SIN CITAR (auditoría A1/A4, 2026-10-19): comentada para que no se imprima; descomentar si se cita.
%   \bibitem{Radhi2025QuantumDeepLearning}
%     Radhi, E. A., Kamil, M. Y., \& Mohammed, M. A. (2025). Quantum machine and deep learning for medical image
%     classification: A systematic review of trends, methodologies, and future directions. \textit{Iraqi Journal for Computer
%     Science and Mathematics}, 6(2), Article 9.
%     \url{https://doi.org/10.52866/2788-7421.1252}

    \bibitem{nielsen2010quantum}
    Nielsen, M. A., \& Chuang, I. L. (2010). \textit{Quantum computation and quantum information: 10th anniversary edition}.
    Cambridge University Press. \url{https://doi.org/10.1017/CBO9780511976667}

  % --- Referencia 31: Variational Quantum Algorithms ---
  \bibitem{cerezo2021variational}
    Cerezo, M., Arrasmith, A., Babbush, R., Benjamin, S. C., Endo, S., Fujii, K., McClean, J. R., Mitarai, K., Yuan, X., Cincio, L., \& Coles, P. J. (2021). Variational quantum algorithms. \textit{Nature Reviews Physics}, 3, 625--644. \url{https://doi.org/10.1038/s42254-021-00348-9}

  % --- Referencia 32: Quantum Circuit Learning ---
  \bibitem{mitarai2018quantum}
    Mitarai, K., Negoro, M., Kitagawa, M., \& Fujii, K. (2018). Quantum circuit learning. \textit{Physical Review A}, 98(3), 032309. \url{https://doi.org/10.1103/PhysRevA.98.032309}

  % --- Referencia 33: Qiskit Circuit Library Documentation ---
  \bibitem{ibmquantum2024docs}
    IBM Quantum. (2024). Qiskit Circuit Library: \texttt{zz\_feature\_map} and \texttt{pauli\_feature\_map}. IBM Quantum Documentation. \url{https://docs.quantum.ibm.com/api/qiskit/qiskit.circuit.library.zz_feature_map}

  % --- Referencia 34: Qiskit Machine Learning Documentation ---
  \bibitem{qiskitml2024}
    IBM Quantum. (2024). Qiskit Machine Learning: \texttt{EstimatorQNN} and \texttt{TorchConnector}. Qiskit Machine Learning Documentation. \url{https://qiskit-community.github.io/qiskit-machine-learning/}

  % --- Referencia 35: Deep Learning ---
  \bibitem{goodfellow2016deep}
    Goodfellow, I., Bengio, Y., \& Courville, A. (2016). \textit{Deep Learning}. MIT Press.
    \url{https://www.deeplearningbook.org}

  % --- Referencia 36: Kernel-Target Alignment ---
  \bibitem{cristianini2002kernel}
    Cristianini, N., Shawe-Taylor, J., Elisseeff, A., \& Kandola, J. (2002). On kernel-target alignment.
    In \textit{Advances in Neural Information Processing Systems (NeurIPS)}, 14.
    \url{https://proceedings.neurips.cc/paper/2001/hash/1f71e393b3809197ed66df836fe833e5-Abstract.html}

  % --- Referencia 37: Davies-Bouldin Cluster Separation Measure ---
  \bibitem{davies1979cluster}
    Davies, D. L., \& Bouldin, D. W. (1979). A cluster separation measure.
    \textit{IEEE Transactions on Pattern Analysis and Machine Intelligence}, 1(2), 224--227.
    \url{https://doi.org/10.1109/TPAMI.1979.4766909}

  % --- Referencia 38: Fisher’s Linear Discriminant ---
  \bibitem{fisher1936use}
    Fisher, R. A. (1936). The use of multiple measurements in taxonomic problems.
    \textit{Annals of Eugenics}, 7(2), 179--188.
    \url{https://doi.org/10.1111/j.1469-1809.1936.tb02137.x}

  % --- Referencia 39: t-SNE Data Visualization ---
  \bibitem{vandermaaten2008tsne}
    van der Maaten, L., \& Hinton, G. (2008). Visualizing data using \textit{t-SNE}.
    \textit{Journal of Machine Learning Research}, 9, 2579--2605.
    \url{https://jmlr.org/papers/v9/vandermaaten08a.html}

  % --- Referencia 40: Evaluation Metrics for Machine Learning ---
  % SIN CITAR (auditoría A1/A4, 2026-10-19): comentada para que no se imprima; descomentar si se cita.
%   \bibitem{powers2011evaluation}
%     Powers, D. M. W. (2011). Evaluation: from precision, recall and \textit{F-measure} to \textit{ROC}, informedness, markedness and correlation.
%     \textit{Journal of Machine Learning Technologies}, 2(1), 37--63.
%     \url{https://arxiv.org/abs/2010.16061}

  % --- Referencia 41: Quantum Machine Learning in Feature Hilbert Spaces ---
  % RETIRADA 2026-10-19 (auditoría A1): duplicado de Schuld2019QuantumML (mismo DOI)
%   \bibitem{schuld2019quantum}
%     Schuld, M., \& Killoran, N. (2019). Quantum machine learning in feature {H}ilbert spaces.
%     \textit{Physical Review Letters}, 122(4), 040504.
%     \url{https://doi.org/10.1103/PhysRevLett.122.040504}

  % --- Referencia 42: Supervised Learning with Quantum-Enhanced Feature Spaces ---
  \bibitem{havlicek2019supervised}
    Havl{\'{i}}{\v{c}}ek, V., C{\'{o}}rcoles, A. D., Temme, K., Harrow, A. W., Kandala, A., Chow, J. M., \& Gambetta, J. M. (2019). Supervised learning with quantum-enhanced feature spaces.
    \textit{Nature}, 567, 209--212.
    \url{https://doi.org/10.1038/s41586-019-0980-2}

  % --- Referencia 43: Power of Data in Quantum Machine Learning ---
  \bibitem{huang2021power}
    Huang, H.-Y., Broughton, M., Mohseni, M., Babbush, R., Boixo, S., Neven, H., \& McClean, J. R. (2021). Power of data in quantum machine learning.
    \textit{Nature Communications}, 12, 2631.
    \url{https://doi.org/10.1038/s41467-021-22539-9}

  % --- Referencia 44: Realization of Contrast Limited Adaptive Histogram Equalization ---
  \bibitem{reza2004realization}
    Reza, A. M. (2004). Realization of the contrast limited adaptive histogram equalization (CLAHE) for real-time image enhancement.
    \textit{Journal of VLSI Signal Processing}, 38, 35--44.
    \url{https://doi.org/10.1023/B:VLSI.0000028532.53893.82}

  % --- Referencia 45: Computational Radiomics System to Decode the Radiographic Phenotype ---
  \bibitem{vanGriethuysen2017}
    van Griethuysen, J. J. M., et al. (2017). Computational radiomics system to decode the radiographic phenotype.
    \textit{Cancer Research}, 77(21), e104--e107.
    \url{https://doi.org/10.1158/0008-5472.CAN-17-0339}

  % --- Referencia 46: IBM Quantum Learning - Basics of Quantum Information ---
  % SIN CITAR (auditoría A1/A4, 2026-10-19): comentada para que no se imprima; descomentar si se cita.
%   \bibitem{ibmquantum2024learning}
%   IBM Quantum. (2024). Basics of quantum information.
%   IBM Quantum Learning.
%   \url{https://learning.quantum.ibm.com/course/basics-of-quantum-information}

\bibitem{Bray2024GLOBOCAN}
    Bray, F., Laversanne, M., Sung, H., Ferlay, J., Siegel, R. L., Soerjomataram, I., \& Jemal, A. (2024).
    Global cancer statistics 2022: GLOBOCAN estimates of incidence and mortality worldwide for 36 cancers in 185 countries.
    \textit{CA: A Cancer Journal for Clinicians}, 74(3), 229--263.
    \url{https://doi.org/10.3322/caac.21834}

  

  \bibitem{Cortes2012}
    Cortes, C., Mohri, M., \& Rostamizadeh, A. (2012).
    Algorithms for learning kernels based on centered alignment.
    \textit{Journal of Machine Learning Research}, 13(28), 795--828.
    \url{https://jmlr.org/papers/v13/cortes12a.html}
  % RETIRADA 2026-10-19 (auditoría A1): duplicado de cristianini2002kernel, sin citar
% 
%   \bibitem{Cristianini2001}
%     Cristianini, N., Shawe-Taylor, J., Elisseeff, A., \& Kandola, J. (2001).
%     On kernel-target alignment. In T. Dietterich, S. Becker, \& Z. Ghahramani (Eds.),
%     \textit{Advances in Neural Information Processing Systems} (Vol. 14, pp. 367--373). MIT Press.
%     \url{https://proceedings.neurips.cc/paper/2001/hash/1f71e393b3809197ed66df836fe833e5-Abstract.html}

  \bibitem{Thanasilp2024}
    Thanasilp, S., Wang, S., Cerezo, M., \& Holmes, Z. (2024).
    Exponential concentration in quantum kernel methods.
    \textit{Nature Communications}, 15(1), 5200.
    \url{https://doi.org/10.1038/s41467-024-49287-w}

  \bibitem{Schuld2021Encoding}
    Schuld, M., Sweke, R., \& Meyer, J. J. (2021).
    Effect of data encoding on the expressive power of variational quantum-machine-learning models.
    \textit{Physical Review A}, 103(3), 032430.
    \url{https://doi.org/10.1103/PhysRevA.103.032430}

  \bibitem{Liu2021Speedup}
    Liu, Y., Arunachalam, S., \& Temme, K. (2021).
    A rigorous and robust quantum speed-up in supervised machine learning.
    \textit{Nature Physics}, 17(9), 1013--1017.
    \url{https://doi.org/10.1038/s41567-021-01287-z}

  \bibitem{Glick2024Covariant}
    Glick, J. R., Gujarati, T. P., C{\'{o}}rcoles, A. D., Kim, Y., Kandala, A., Gambetta, J. M., \& Temme, K. (2024).
    Covariant quantum kernels for data with group structure.
    \textit{Nature Physics}, 20(3), 479--483.
    \url{https://doi.org/10.1038/s41567-023-02340-9}

  \bibitem{Huang2022Experiments}
    Huang, H.-Y., Broughton, M., Cotler, J., Chen, S., Li, J., Mohseni, M., Neven, H., Babbush, R., Kueng, R., Preskill, J., \& McClean, J. R. (2022).
    Quantum advantage in learning from experiments.
    \textit{Science}, 376(6598), 1182--1186.
    \url{https://doi.org/10.1126/science.abn7293}

  \bibitem{Hubregtsen2022}
    Hubregtsen, T., Wierichs, D., Gil-Fuster, E., Derks, P.-J. H. S., Faehrmann, P. K., \& Meyer, J. J. (2022).
    Training quantum embedding kernels on near-term quantum computers.
    \textit{Physical Review A}, 106(4), 042431.
    \url{https://doi.org/10.1103/PhysRevA.106.042431}

  \bibitem{Guyon2003}
    Guyon, I., \& Elisseeff, A. (2003).
    An introduction to variable and feature selection.
    \textit{Journal of Machine Learning Research}, 3, 1157--1182.
    \url{https://jmlr.org/papers/v3/guyon03a.html}

  \bibitem{ElHafidi2025}
    Toufah, A., Kadim, M. A., \& El Hafidi, M. Y. (2025).
    Investigating quantum feature maps in quantum support vector machines for lung cancer classification.
    \textit{arXiv preprint} arXiv:2506.03272.
    \url{https://doi.org/10.48550/arXiv.2506.03272}

  \bibitem{Arrasmith2021}
    Arrasmith, A., Cerezo, M., Czarnik, P., Cincio, L., \& Coles, P. J. (2021).
    Effect of barren plateaus on gradient-free optimization.
    \textit{Quantum}, 5, 558.
    \url{https://doi.org/10.22331/q-2021-10-05-558}

  \bibitem{Rifat2024}
    Rifat, R. C., Tonmoy, M. T. U. I., Aziz, R. T., \& Abdul Qayum, M. (2024).
    Enhanced hybrid quantum neural network for breast cancer detection.
    En \textit{2024 27th International Conference on Computer and Information Technology (ICCIT)} (pp. 1464--1469). IEEE.
    \url{https://doi.org/10.1109/iccit64611.2024.11021989}

  \bibitem{Yang2023MedMNIST}
    Yang, J., Shi, R., Wei, D., Liu, Z., Zhao, L., Ke, B., Pfister, H., \& Ni, B. (2023).
    MedMNIST v2 -- A large-scale lightweight benchmark for 2D and 3D biomedical image classification.
    \textit{Scientific Data}, 10(1), 41.
    \url{https://doi.org/10.1038/s41597-022-01721-8}

  \bibitem{Gillies2016Radiomics}
    Gillies, R. J., Kinahan, P. E., \& Hricak, H. (2016).
    Radiomics: Images are more than pictures, they are data.
    \textit{Radiology}, 278(2), 563--577.
    \url{https://doi.org/10.1148/radiol.2015151169}

  \bibitem{Scholkopf1998KPCA}
    Sch{\"o}lkopf, B., Smola, A., \& M{\"u}ller, K.-R. (1998).
    Nonlinear component analysis as a kernel eigenvalue problem.
    \textit{Neural Computation}, 10(5), 1299--1319.
    \url{https://doi.org/10.1162/089976698300017467}

  \bibitem{CortesVapnik1995}
    Cortes, C., \& Vapnik, V. (1995).
    Support-vector networks.
    \textit{Machine Learning}, 20(3), 273--297.
    \url{https://doi.org/10.1007/BF00994018}

\bibitem{Lee2017CBISDDSM}
    Lee, R. S., Gimenez, F., Hoogi, A., Miyake, K. K., Gorovoy, M., \& Rubin, D. L. (2017).
    A curated mammography data set for use in computer-aided detection and diagnosis research.
    \textit{Scientific Data}, 4(1), 170177.
    \url{https://doi.org/10.1038/sdata.2017.177}

  \bibitem{Clark2013TCIA}
    Clark, K., Vendt, B., Smith, K., Freymann, J., Kirby, J., Koppel, P., Moore, S., Phillips, S., Maffitt, D., Pringle, M., Tarbox, L., \& Prior, F. (2013).
    The Cancer Imaging Archive (TCIA): Maintaining and operating a public information repository.
    \textit{Journal of Digital Imaging}, 26(6), 1045--1057.
    \url{https://doi.org/10.1007/s10278-013-9622-7}

  \bibitem{Zwanenburg2020IBSI}
    Zwanenburg, A., Valli\`eres, M., Abdalah, M. A., Aerts, H. J. W. L., Andrearczyk, V., Apte, A., Ashrafinia, S., Bakas, S., Beukinga, R. J., Boellaard, R., Bogowicz, M., Boldrini, L., Buvat, I., Cook, G. J. R., Davatzikos, C., Depeursinge, A., Desseroit, M.-C., Dinapoli, N., Dinh, C. V., \ldots\ L\"ock, S. (2020).
    The Image Biomarker Standardization Initiative: Standardized quantitative radiomics for high-throughput image-based phenotyping.
    \textit{Radiology}, 295(2), 328--338.
    \url{https://doi.org/10.1148/radiol.2020191145}

  \bibitem{Preskill2018NISQ}
    Preskill, J. (2018).
    Quantum computing in the NISQ era and beyond.
    \textit{Quantum}, 2, 79.
    \url{https://doi.org/10.22331/q-2018-08-06-79}

  \bibitem{McClean2018Barren}
    McClean, J. R., Boixo, S., Smelyanskiy, V. N., Babbush, R., \& Neven, H. (2018).
    Barren plateaus in quantum neural network training landscapes.
    \textit{Nature Communications}, 9(1), 4812.
    \url{https://doi.org/10.1038/s41467-018-07090-4}

  \bibitem{Mari2020Transfer}
    Mari, A., Bromley, T. R., Izaac, J., Schuld, M., \& Killoran, N. (2020).
    Transfer learning in hybrid classical-quantum neural networks.
    \textit{Quantum}, 4, 340.
    \url{https://doi.org/10.22331/q-2020-10-09-340}

  \bibitem{SchuldKilloran2022Advantage}
    Schuld, M., \& Killoran, N. (2022).
    Is quantum advantage the right goal for quantum machine learning?
    \textit{PRX Quantum}, 3(3), 030101.
    \url{https://doi.org/10.1103/PRXQuantum.3.030101}

  \bibitem{Shaydulin2022Bandwidth}
    Shaydulin, R., \& Wild, S. M. (2022).
    Importance of kernel bandwidth in quantum machine learning.
    \textit{Physical Review A}, 106(4), 042407.
    \url{https://doi.org/10.1103/PhysRevA.106.042407}

  \bibitem{Canatar2023Bandwidth}
    Canatar, A., Peters, E., Pehlevan, C., Wild, S. M., \& Shaydulin, R. (2022).
    Bandwidth enables generalization in quantum kernel models.
    \textit{arXiv preprint} arXiv:2206.06686.
    \url{https://doi.org/10.48550/arXiv.2206.06686}

  \bibitem{Kubler2021Inductive}
    K\"ubler, J. M., Buchholz, S., \& Sch\"olkopf, B. (2021).
    The inductive bias of quantum kernels.
    En \textit{Advances in Neural Information Processing Systems 34 (NeurIPS 2021)}.
    \url{https://doi.org/10.48550/arXiv.2106.03747}

  \bibitem{Bowles2024Benchmarking}
    Bowles, J., Ahmed, S., \& Schuld, M. (2024).
    Better than classical? The subtle art of benchmarking quantum machine learning models.
    \textit{arXiv preprint} arXiv:2403.07059.
    \url{https://doi.org/10.48550/arXiv.2403.07059}

  \bibitem{Le2011FRQI}
    Le, P. Q., Dong, F., \& Hirota, K. (2011).
    A flexible representation of quantum images for polynomial preparation, image compression, and processing operations.
    \textit{Quantum Information Processing}, 10(1), 63--84.
    \url{https://doi.org/10.1007/s11128-010-0177-y}

  \bibitem{Zhang2013NEQR}
    Zhang, Y., Lu, K., Gao, Y., \& Wang, M. (2013).
    NEQR: A novel enhanced quantum representation of digital images.
    \textit{Quantum Information Processing}, 12(8), 2833--2860.
    \url{https://doi.org/10.1007/s11128-013-0567-z}

  \bibitem{Parigi2026QImR}
    Parigi, M., Khosrojerdi, M., Caruso, F., \& Banchi, L. (2026).
    Analysis of quantum image representations for supervised classification.
    \textit{AVS Quantum Science}, 8(1), 013801.
    \url{https://doi.org/10.1116/5.0296376}


  \bibitem{Elmore1994}
    Elmore, J. G., Wells, C. K., Lee, C. H., Howard, D. H., \& Feinstein, A. R. (1994).
    Variability in radiologists' interpretations of mammograms.
    \textit{New England Journal of Medicine}, 331(22), 1493--1499.
    \url{https://doi.org/10.1056/NEJM199412013312206}

  \bibitem{Monticciolo2021ACR}
    Monticciolo, D. L., Malak, S. F., Friedewald, S. M., Eby, P. R., Newell, M. S., Moy, L.,
    Destounis, S., Leung, J. W. T., Hendrick, R. E., \& Smetherman, D. (2021).
    Breast cancer screening recommendations inclusive of all women at average risk: Update from the ACR and Society of Breast Imaging.
    \textit{Journal of the American College of Radiology}, 18(9), 1280--1288.
    \url{https://doi.org/10.1016/j.jacr.2021.04.021}
\end{thebibliography}
